%% Elsevier submission version -- Social Networks
%% Compile with: pdflatex elsevier -> pdflatex elsevier
%% The EAI-formatted source (icst.cls) is retained as main.tex.
\documentclass[preprint,12pt,authoryear]{elsarticle}

\usepackage{amsmath}
\usepackage{amssymb}
\usepackage{amsthm}
\usepackage{booktabs}
\usepackage{graphicx}
\usepackage{multirow}
\usepackage{xcolor}
\usepackage{url}
\usepackage{subfiles}
\usepackage[breaklinks,colorlinks,bookmarksopen,linkcolor=blue,citecolor=blue,
            urlcolor=blue]{hyperref}

\makeatletter
\g@addto@macro{\UrlBreaks}{\do\-}
\makeatother

\newtheorem{theorem}{Theorem}
\newtheorem{proposition}[theorem]{Proposition}
\newtheorem{lemma}[theorem]{Lemma}
\newtheorem{corollary}[theorem]{Corollary}
\theoremstyle{definition}
\newtheorem{definition}[theorem]{Definition}
\newtheorem{assumption}[theorem]{Assumption}
\newtheorem{remark}[theorem]{Remark}

%% elsarticle has no \email macro of its own in this context
\providecommand{\email}[1]{\texttt{#1}}

\journal{Social Networks}

\begin{document}

%% Anonymised for double-anonymized review.
%% Author information is in the separate title-page file.

\begin{frontmatter}

\title{The Camouflage Ratio: Strategic Tie Formation and the\\
       Limits of Structural Detection}

\begin{abstract}

Inferring an actor's attributes from those of its network neighbours is among the
oldest uses of relational data, and it rests on homophily: neighbours are assumed to
be informative. We ask what happens when an actor \emph{chooses} its ties in order to
defeat that inference. Automated accounts on social platforms do exactly this,
forming ties to genuine users so that their local neighbourhood resembles a genuine
one---a strategy known as relation camouflage. We formalise it with a single
structural parameter, the \emph{camouflage ratio} $\rho$, the share of an actor's
ties directed at the other class, and ask when neighbourhood information still
improves on attributes alone.

Three results follow. First, camouflage is not a private act: because a tie has two
ends, an actor that camouflages necessarily contaminates its neighbours'
neighbourhoods in turn. Edge-endpoint conservation forces the induced contamination
of the majority to be $\eta = c\rho$, where $c$ is the minority-to-majority ratio, so
that \emph{prevalence acts as protection}---the more common the strategic actors, the
less camouflage each one needs to defeat relational inference. Second, this yields a
threshold $\rho^\star = (1-D^{-1/2})/(1+c)$ in the mean neighbourhood size $D$, above
which averaging over neighbours is worse than ignoring the network entirely, and---in
the noise-dominated regime, and only where the minority is large enough relative to
$D^{-1/2}$---a bounded interval above which relational information becomes informative
again with its sign reversed, so that within the model the most damaging strategy is
\emph{intermediate} rather than maximal camouflage. We show that the networks we
measure do not fall in that regime, and are harmful without any such reversal. We are explicit that a crossover of this kind is not new---at $c = 1$ it
recovers a result of Ma et al.\ (2022) and the interval is the second root of the
same condition---and that our contribution is the asymmetric, class-imbalanced case
together with a variance correction that shifts the threshold by $11.5\%$. Third,
selectively discarding suspicious ties does not straightforwardly restore inference:
it trades bias against the number of ties left to average over, and we bound the
discrimination achievable from attribute discrepancies alone.

We then measure $\rho$ on MGTAB, a Twitter dataset of $10{,}199$ expert-annotated
accounts in which every account is labelled, so that no tie is lost to unlabelled
alters. Evaluating the exact criterion at each relation's own parameters, \emph{every}
recorded relation turns out to be one on which averaging over alters is worse than
using attributes alone, by factors between $1.3$ and $3.6$. The layers nevertheless
differ sharply in mechanism: on every social relation the automated accounts are almost
perfectly camouflaged---the median one has \emph{no} automated alter, and
automated--automated ties number $212$ against $40{,}263$ automated--genuine ties---
whereas on content-coordination relations (shared URLs, shared hashtags) the same
accounts converge, with automated--automated ties rising to $15{,}963$ and $28{,}203$.
They avoid one another where ties are visible and meet where content is shared, and
$\rho$ differs between the layers by a factor approaching two on identical actors, so
multiplexity rather than aggregate density decides whether relational inference is
usable. We also document why the canonical bot benchmarks cannot support this
measurement at all, and confirm the edge-endpoint identity to machine precision in its
degree-corrected form.

\end{abstract}

\begin{keyword}
Relational inference \sep Homophily \sep Strategic tie formation \sep Multiplexity \sep Social bots \sep Stochastic block models
\end{keyword}

\end{frontmatter}

\subfile{sections/1_intro}
\subfile{sections/2_related}
\subfile{sections/3_model}
\subfile{sections/4_theory}
\subfile{sections/5_validation}
\subfile{sections/6_discussion}

\section*{Declaration of competing interest}

The authors declare that they have no known competing financial interests or personal
relationships that could have appeared to influence the work reported in this paper.

\section*{Data availability}

All data analysed in this study are third-party datasets that are publicly available
and were not generated by the authors. The Twitter account data are from MGTAB
\citep{shi2023mgtab}, distributed by its authors under a Creative Commons
Attribution--NonCommercial--NoDerivatives 4.0 licence; we analysed the released
tensors without modification and redistribute no derivative of them. The two review
networks are those of \citet{dou2020caregnn} and are openly downloadable. The
cresci-2015 and cresci-2017 collections referred to in
Section~\ref{sec:measuring} are likewise public. TwiBot-20 and TwiBot-22 require
permission from their authors, which we did not obtain and which is why they are not
analysed here.

\begin{thebibliography}{99}

\bibitem[Ajorlou et al.(2025)]{ajorlou2025dirichlet}
Ajorlou, H., Mateos, G., Ruiz, L., 2025. Dirichlet meets Horvitz and Thompson: estimating homophily in large networks via sampling. In: 2025 59th Asilomar Conference on Signals, Systems, and Computers. IEEE, pp. 1109-1113. doi:10.1109/IEEECONF67917.2025.11443313

\bibitem[An et al.(2026)]{an2026estimating}
An, W., Estrada, P., Estrada, J., Jacho-Chavez, D., 2026. Estimating peer influence in multilayer networks. Soc Networks.86:252-65. doi:10.1016/j.socnet.2026.03.003

\bibitem[Anderson et al.(1999)]{anderson1999interaction}
Anderson, B.S., Butts, C., Carley, K., 1999. The interaction of size and density with graph-level indices. Soc Networks.21(3):239-67. doi:10.1016/S0378-8733(99)00011-8

\bibitem[Ariyarathne et al.(2026)]{ariyarathne2026behavior}
Ariyarathne, I., Ariyarathne, G., Flammini, A., Menczer, F., Nwala, A., 2026. Behavior change as a signal for identifying social media manipulation. In: Proceedings of the 18th ACM Web Science Conference (WebSci '26). ACM, pp. 349-358. doi:10.1145/3795766.3799748

\bibitem[Baranwal et al.(2021)]{baranwal2021graph}
Baranwal, A., Fountoulakis, K., Jagannath, A., 2021. Graph convolution for semi-supervised classification: improved linear separability and out-of-distribution generalization. In: Proceedings of the 38th International Conference on Machine Learning (ICML 2021). PMLR.139:684-93.

\bibitem[Baranwal et al.(2023a)]{baranwal2023effects}
Baranwal, A., Fountoulakis, K., Jagannath, A., 2023a. Effects of graph convolutions in multi-layer networks. In: Proceedings of the 11th International Conference on Learning Representations (ICLR 2023).

\bibitem[Baranwal et al.(2023b)]{baranwal2023optimality}
Baranwal, A., Fountoulakis, K., Jagannath, A., 2023b. Optimality of message-passing architectures for sparse graphs. In: Advances in Neural Information Processing Systems 36 (NeurIPS 2023).

\bibitem[Baranwal(2026)]{baranwal2026depth}
Baranwal, A.R., 2026. The value of depth in message passing on sparse graphs: a Kesten--Stigum dichotomy. arXiv preprint arXiv:2607.16676.

\bibitem[Barone and Coscia(2018)]{barone2018birds}
Barone, M., Coscia, M., 2018. Birds of a feather scam together: trustworthiness homophily in a business network. Soc Networks.54:228-37. doi:10.1016/j.socnet.2018.01.009

\bibitem[Blanchard et al.(2017)]{blanchard2017machine}
Blanchard, P., El Mhamdi, E.M., Guerraoui, R., Stainer, J., 2017. Machine learning with adversaries: Byzantine tolerant gradient descent. In: Advances in Neural Information Processing Systems 30 (NeurIPS 2017). p. 118-28.

\bibitem[Bojanowski and Corten(2014)]{bojanowski2014measuring}
Bojanowski, M., Corten, R., 2014. Measuring segregation in social networks. Soc Networks.39:14-32. doi:10.1016/j.socnet.2014.04.001

\bibitem[Bojchevski and G\"unnemann(2019)]{bojchevski2019certifiable}
Bojchevski, A., G\"unnemann, S., 2019. Certifiable robustness to graph perturbations. In: Advances in Neural Information Processing Systems 32 (NeurIPS 2019).

\bibitem[Burt et al.(2021)]{burt2021network}
Burt, R.S., Reagans, R.E., Volvovsky, H.C., 2021. Network brokerage and the perception of leadership. Soc Networks.65:33-50. doi:10.1016/j.socnet.2020.09.002

\bibitem[Chandrasekhar et al.(2026)]{chandrasekhar2026multiplexing}
Chandrasekhar, A.G., Chaudhary, V., Golub, B., Jackson, M.O., 2026. Multiplexing in networks and diffusion. Proc. Natl. Acad. Sci. U.S.A. 123, e2534923123. doi:10.1073/pnas.2534923123

\bibitem[Chen et al.(2021)]{chen2021understanding}
Chen, L., Li, J., Peng, Q., Liu, Y., Zheng, Z., Yang, C., 2021. Understanding structural vulnerability in graph convolutional networks. In: Proceedings of the 30th International Joint Conference on Artificial Intelligence (IJCAI 2021). p. 2249-55. doi:10.24963/ijcai.2021/310

\bibitem[Choi and Young(2026)]{choi2026who}
Choi, E.C., Young, L.E., 2026. Who you know, what you believe, what you share: personal networks and health misinformation susceptibility. Soc Networks.87:163-72. doi:10.1016/j.socnet.2026.06.008

\bibitem[Cresci(2020)]{cresci2020decade}
Cresci, S., 2020. A decade of social bot detection. Commun ACM.63(10):72-83. doi:10.1145/3409116

\bibitem[Dang et al.(2026)]{dang2026xhbot}
Dang, Q.-V., Nguyen, P.-L., Le, D., Dinh, M.N., 2026. XHBot: eXplainable heterophily-aware graph neural networks for social bot detection. EAI Endorsed Trans AI Robotics [Internet]. 2026 Jul 7 [cited 2026 Jul 31];5. Available from: \url{https://publications.eai.eu/index.php/airo/article/view/12969}

\bibitem[Deshpande et al.(2018)]{deshpande2018contextual}
Deshpande, Y., Sen, S., Montanari, A., Mossel, E., 2018. Contextual stochastic block models. In: Advances in Neural Information Processing Systems 31 (NeurIPS 2018). p. 8590-602.

\bibitem[Dou et al.(2020)]{dou2020caregnn}
Dou, Y., Liu, Z., Sun, L., Deng, Y., Peng, H., Yu, P.S., 2020. Enhancing graph neural network-based fraud detectors against camouflaged fraudsters. In: Proceedings of the 29th ACM International Conference on Information and Knowledge Management (CIKM 2020). ACM. doi:10.1145/3340531.3411903

\bibitem[Duranthon and Zdeborov\'a(2025)]{duranthon2025statistical}
Duranthon, O., Zdeborov\'a, L., 2025. Statistical physics analysis of graph neural networks: approaching optimality in the contextual stochastic block model. Phys. Rev. X 15, 041026. doi:10.1103/lfxj-hbsk

\bibitem[Espinosa-Rada et al.(2026)]{espinosarada2026socioeconomic}
Espinosa-Rada, A., Bargsted, M., Ortiz Ruiz, F., 2026. Socioeconomic homogeneity in acquaintance networks: occupational prestige, education, and intergenerational mobility. Soc. Networks 86, 342-353. doi:10.1016/j.socnet.2026.04.012

\bibitem[Everett and Borgatti(2026)]{everett2026alter}
Everett, M.G., Borgatti, S.P., 2026. Alter composition with overlapping group memberships. Soc Networks.85:80-8. doi:10.1016/j.socnet.2025.12.001

\bibitem[Fluer et al.(2026)]{fluer2026survey}
Fluer, A., Laga, I., Graham, L., Almirol, E., Meyer, M., Schneider, J.A., Cummins, B., 2026. From survey data to social multiplex models: incorporating interlayer correlation from multiple data sources. Soc. Networks 86, 104-119. doi:10.1016/j.socnet.2026.01.005

\bibitem[Fountoulakis et al.(2023)]{fountoulakis2023graph}
Fountoulakis, K., Levi, A., Yang, S., Baranwal, A., Jagannath, A., 2023. Graph attention retrospective. J Mach Learn Res.24(246):1-52. Available from: \url{https://jmlr.org/papers/volume24/22-125/22-125.pdf}

\bibitem[Geisler et al.(2020)]{geisler2020reliable}
Geisler, S., Z\"ugner, D., G\"unnemann, S., 2020. Reliable graph neural networks via robust aggregation. In: Advances in Neural Information Processing Systems 33 (NeurIPS 2020).

\bibitem[Gosch et al.(2023)]{gosch2023revisiting}
Gosch, L., Geisler, S., Sturm, D., Charpentier, B., Z\"ugner, D., G\"unnemann, S., 2023. Revisiting robustness in graph machine learning. In: Proceedings of the 11th International Conference on Learning Representations (ICLR 2023).

\bibitem[He et al.(2025)]{he2025bothp}
He, B., Jiang, X., Wu, Q., Liu, H., Yang, Y., Liao, Y., 2025. Boosting bot detection via heterophily-aware representation learning and prototype-guided cluster discovery. In: Proceedings of the 31st ACM SIGKDD Conference on Knowledge Discovery and Data Mining (KDD 2025). ACM. doi:10.1145/3711896.3736862

\bibitem[Holland et al.(1983)]{holland1983stochastic}
Holland, P.W., Laskey, K.B., Leinhardt, S., 1983. Stochastic blockmodels: first steps. Soc Networks.5(2):109-37. doi:10.1016/0378-8733(83)90021-7

\bibitem[Jin et al.(2020)]{jin2020graph}
Jin, W., Ma, Y., Liu, X., Tang, X., Wang, S., Tang, J., 2020. Graph structure learning for robust graph neural networks. In: Proceedings of the 26th ACM SIGKDD International Conference on Knowledge Discovery and Data Mining (KDD 2020). ACM. p. 66-74. doi:10.1145/3394486.3403049

\bibitem[Kami\'nski et al.(2026)]{kaminski2026relationships}
Kami\'nski, B., Pra{\l}at, P., Wojnarowicz, A., Zawisza, M., 2026. Relationships between node degrees and hyperedge sizes in empirical hypergraphs. Soc. Networks 87, 119-134. doi:10.1016/j.socnet.2026.06.004

\bibitem[Krivitsky et al.(2009)]{krivitsky2009representing}
Krivitsky, P.N., Handcock, M.S., Raftery, A.E., Hoff, P.D., 2009. Representing degree distributions, clustering, and homophily in social networks with latent cluster random effects models. Soc Networks.31(3):204-13. doi:10.1016/j.socnet.2009.04.001

\bibitem[Lakdawala et al.(2026)]{lakdawala2026bonds}
Lakdawala, R., Leenders, R., Mulder, J., 2026. Not all bonds are created equal: dyadic latent class models for relational event data. Soc. Networks 87, 135-151. doi:10.1016/j.socnet.2026.06.006

\bibitem[Lazega and Pattison(1999)]{lazega1999multiplexity}
Lazega, E., Pattison, P.E., 1999. Multiplexity, generalized exchange and cooperation in organizations: a case study. Soc Networks.21(1):67-90. doi:10.1016/S0378-8733(99)00002-7

\bibitem[Leenders(2002)]{leenders2002modeling}
Leenders, R.T.h.A.J., 2002. Modeling social influence through network autocorrelation: constructing the weight matrix. Soc Networks.24(1):21-47. doi:10.1016/S0378-8733(01)00049-1

\bibitem[Li et al.(2018)]{li2018deeper}
Li, Q., Han, Z., Wu, X.M., 2018. Deeper insights into graph convolutional networks for semi-supervised learning. In: Proceedings of the 32nd AAAI Conference on Artificial Intelligence (AAAI 2018). p. 3538-45. doi:10.1609/aaai.v32i1.11604

\bibitem[Luan et al.(2022)]{luan2022revisiting}
Luan, S., Hua, C., Lu, Q., Zhu, J., Zhao, M., Zhang, S., et al., 2022. Revisiting heterophily for graph neural networks. In: Advances in Neural Information Processing Systems 35 (NeurIPS 2022).

\bibitem[Luan et al.(2023)]{luan2023when}
Luan, S., Hua, C., Lu, Q., Zhu, J., Zhao, M., Zhang, S., et al., 2023. When do graph neural networks help with node classification? Investigating the impact of homophily principle on node distinguishability. In: Advances in Neural Information Processing Systems 36 (NeurIPS 2023).

\bibitem[Ma et al.(2022)]{ma2022homophily}
Ma, Y., Liu, X., Shah, N., Tang, J., 2022. Is homophily a necessity for graph neural networks? In: International Conference on Learning Representations (ICLR 2022).

\bibitem[Ma et al.(2025)]{ma2025attention}
Ma, Z., Zhang, Q., Zhou, B., Zhang, Y., Hu, S., Wang, Z., 2025. Graph attention is not always beneficial: a theoretical analysis of graph attention mechanisms via contextual stochastic block models. In: Proceedings of the 42nd International Conference on Machine Learning (ICML 2025). arXiv:2412.15496.

\bibitem[Macdonald(2026)]{macdonald2026wiretap}
Macdonald, M., 2026. Wiretap surveillance of criminal networks does not resemble other types of collaborative networks. Soc. Networks 87, 12-23. doi:10.1016/j.socnet.2026.05.002

\bibitem[Mao et al.(2023)]{mao2023demystifying}
Mao, H., Chen, Z., Jin, W., Han, H., Ma, Y., Zhao, T., et al., 2023. Demystifying structural disparity in graph neural networks: can one size fit all? In: Advances in Neural Information Processing Systems 36 (NeurIPS 2023).

\bibitem[McMillan(2022)]{mcmillan2022worth}
McMillan, C., 2022. Worth the weight: conceptualizing and measuring strong versus weak tie homophily. Soc Networks.68:139-47. doi:10.1016/j.socnet.2021.06.003

\bibitem[Mujkanovic et al.(2022)]{mujkanovic2022are}
Mujkanovic, F., Geisler, S., G\"unnemann, S., Bojchevski, A., 2022. Are defenses for graph neural networks robust? In: Advances in Neural Information Processing Systems 35 (NeurIPS 2022).

\bibitem[Mukherjee et al.(2026)]{mukherjee2026optimal}
Mukherjee, K., Alom, Z., Ngo, T.G.B., Akcora, C.G., Kantarcioglu, M., 2026. Optimal transport-guided adversarial attacks on graph neural network-based bot detection. arXiv preprint arXiv:2602.00318. To appear, Proceedings of the 43rd International Conference on Machine Learning (ICML 2026).

\bibitem[Nguyen(2026)]{nguyen2026agency}
Nguyen, Q.T., 2026. Agency within structure: strategic behavior and network position in explaining premium buyer access in global buyer--seller networks. Soc. Networks 87, 188-209. doi:10.1016/j.socnet.2026.07.004

\bibitem[Noel and Nyhan(2011)]{noel2011unfriending}
Noel, H., Nyhan, B., 2011. The `unfriending' problem: the consequences of homophily in friendship retention for causal estimates of social influence. Soc Networks.33(3):211-8. doi:10.1016/j.socnet.2011.05.003

\bibitem[NT and Maehara(2019)]{nt2019revisiting}
NT, H., Maehara, T., 2019. Revisiting graph neural networks: all we have is low-pass filters. arXiv preprint. doi:10.48550/arXiv.1905.09550

\bibitem[Oono and Suzuki(2020)]{oono2020graph}
Oono, K., Suzuki, T., 2020. Graph neural networks exponentially lose expressive power for node classification. In: International Conference on Learning Representations (ICLR 2020).

\bibitem[Pan et al.(2026)]{pan2026camera}
Pan, J., Liu, Y., Zheng, Y., Chi, L., Liew, A.W.C., Pan, S., 2026. CAMERA: adapting to semantic camouflage in unsupervised text-attributed graph fraud detection. arXiv preprint arXiv:2605.20032. Accepted, IJCAI 2026.

\bibitem[Ran et al.(2025)]{ran2025identifying}
Ran, Y., Xiao, J., Xu, X.K., 2025. Identifying social bots via heterogeneous motifs based on na\"ive Bayes model. IEEE Trans. Inf. Forensics Secur. 20, 12754-12766. doi:10.1109/TIFS.2025.3635006

\bibitem[Roeling and Nicholls(2020)]{roeling2020imputation}
Roeling, M.P., Nicholls, G.K., 2020. Imputation of attributes in networked data using Bayesian autocorrelation regression models. Soc Networks.62:24-32. doi:10.1016/j.socnet.2020.02.005

\bibitem[Sch\"u{\ss}ler and Fuhse(2025)]{schussler2025duality}
Sch\"u{\ss}ler, A., Fuhse, J.A., 2025. The duality of network ties and attributes. Soc. Networks 83, 186-198. doi:10.1016/j.socnet.2025.07.002

\bibitem[Shea et al.(2019)]{shea2019cheater}
Shea, C.T., Lee, J., Menon, T., Im, D.-K., 2019. Cheater's hide and seek: strategic cognitive network activation during ethical decision making. Soc Networks.58:143-55. doi:10.1016/j.socnet.2019.03.005

\bibitem[Shi et al.(2025)]{shi2023mgtab}
Shi, S., Qiao, K., Liu, Z., Yang, J., Chen, C., Chen, J., et al., 2025. MGTAB: a multi-relational graph-based Twitter account detection benchmark. Neurocomputing.647:130490. doi:10.1016/j.neucom.2025.130490

\bibitem[Snijders et al.(2013)]{snijders2013model}
Snijders, T.A.B., Lomi, A., Torl\'o, V.J., 2013. A model for the multiplex dynamics of two-mode and one-mode networks, with an application to employment preference, friendship, and advice. Soc Networks.35(2):265-76. doi:10.1016/j.socnet.2012.05.005

\bibitem[Snijders(1981)]{snijders1981degree}
Snijders, T.A.B., 1981. The degree variance: an index of graph heterogeneity. Soc Networks.3(3):163-74. doi:10.1016/0378-8733(81)90014-9

\bibitem[Tang et al.(2026)]{tang2026twomode}
Tang, H., Mi, M., Ma, Y., Yang, H., Wang, F., 2026. Two-mode network autoregressive model for network analysis with core nodes. Soc. Networks 86, 75-87. doi:10.1016/j.socnet.2026.01.004

\bibitem[V\"or\"os and Snijders(2017)]{voros2017cluster}
V\"or\"os, A., Snijders, T.A.B., 2017. Cluster analysis of multiplex networks: defining composite network measures. Soc Networks.49:93-112. doi:10.1016/j.socnet.2017.01.002

\bibitem[V\"or\"os et al.(2019)]{voros2019limits}
V\"or\"os, A., Block, P., Boda, Z., 2019. Limits to inferring status from friendship relations. Soc Networks.59:77-97. doi:10.1016/j.socnet.2019.05.007

\bibitem[Weiss(2026)]{weiss2026cost}
Weiss, E., 2026. Cost-sensitive neighborhood aggregation for heterophilous graphs: when does per-edge routing help? arXiv preprint. doi:10.48550/arXiv.2603.24291

\bibitem[Yin et al.(2018)]{yin2018byzantine}
Yin, D., Chen, Y., Kannan, R., Bartlett, P., 2018. Byzantine-robust distributed learning: towards optimal statistical rates. In: Proceedings of the 35th International Conference on Machine Learning (ICML 2018). PMLR 80. p. 5650-9.

\bibitem[Yu et al.(2018)]{yu2018detecting}
Yu, L., Woodall, W.H., Tsui, K.-L., 2018. Detecting node propensity changes in the dynamic degree corrected stochastic block model. Soc Networks.54:209-27. doi:10.1016/j.socnet.2018.03.004

\bibitem[Z\"ugner and G\"unnemann(2019)]{zugner2019adversarial}
Z\"ugner, D., G\"unnemann, S., 2019. Adversarial attacks on graph neural networks via meta learning. In: Proceedings of the 7th International Conference on Learning Representations (ICLR 2019).

\bibitem[Z\"ugner et al.(2018)]{zugner2018adversarial}
Z\"ugner, D., Akbarnejad, A., G\"unnemann, S., 2018. Adversarial attacks on neural networks for graph data. In: Proceedings of the 24th ACM SIGKDD International Conference on Knowledge Discovery and Data Mining (KDD 2018). ACM. p. 2847-56. doi:10.1145/3219819.3220078

\bibitem[Zhang and Zitnik(2020)]{zhang2020gnnguard}
Zhang, X., Zitnik, M., 2020. GNNGuard: defending graph neural networks against adversarial attacks. In: Advances in Neural Information Processing Systems 33 (NeurIPS 2020).

\bibitem[Zhu et al.(2020)]{zhu2020beyond}
Zhu, J., Yan, Y., Zhao, L., Heimann, M., Akoglu, L., Koutra, D., 2020. Beyond homophily in graph neural networks: current limitations and effective designs. In: Advances in Neural Information Processing Systems 33 (NeurIPS 2020).

\bibitem[Zhu et al.(2022)]{zhu2022how}
Zhu, J., Jin, J., Loveland, D., Schaub, M.T., Koutra, D., 2022. How does heterophily impact the robustness of graph neural networks? Theoretical connections and practical implications. In: Proceedings of the 28th ACM SIGKDD Conference on Knowledge Discovery and Data Mining (KDD 2022). ACM. p. 2637-47. doi:10.1145/3534678.3539418

\end{thebibliography}

\subfile{sections/7_appendix}

\end{document}
